
\documentclass[11pt]{article}
\usepackage[a4paper,margin=0.95in]{geometry}
\usepackage{fontspec}
\setmainfont{Noto Serif}
\usepackage{microtype}
\usepackage{amsmath,amssymb,mathtools}
\usepackage{array,longtable,booktabs}
\usepackage{graphicx}
\usepackage{csquotes}
\usepackage{hyperref}
\hypersetup{colorlinks=false,pdfborder={0 0 0}}
\setlength{\parindent}{1.15em}
\setlength{\parskip}{0.18em}
\allowdisplaybreaks
\sloppy
\newcommand{\scn}[1]{\textsc{#1}}
\newcommand{\W}{\mathrm{W}}
\begin{document}

\begin{center}
{\Large\bfseries CORRESPONDENCE BETWEEN LEJEUNE DIRICHLET AND GAUSS\footnote{Of Gauss's replies to the letters of Dirichlet printed below, only the two letters which are printed in Gauss, Werke, vol. II, pp. 514--518, were available. According to a communication from E. Schering, no other replies have been found in Gauss's estate. F.}\par}
\vspace{0.8em}
\end{center}

\section*{I. Lejeune Dirichlet to Gauss}
\begin{center}\emph{Most revered Court Councillor!}\end{center}

I have the honour of most respectfully presenting to Your Excellency my first mathematical attempt, as a token of my most heartfelt veneration and my deepest admiration. With the kind recommendation of Baron von \scn{Humboldt}, I venture to flatter myself with the hope that you will receive and judge this first work of a young German with kindly indulgence. If Your Excellency should not find it wholly unworthy of your attention, I would venture most respectfully to ask permission to write to you from time to time and to ask you for some hints for the guidance of my further scientific efforts. I would regard this permission as the highest happiness, since, with the preference with which I pursue the study of indeterminate analysis, I wish for nothing so earnestly as to see the author of the immortal \emph{Disquisitiones arithmeticae} take some interest in my efforts. Since until now I have chiefly occupied myself with higher arithmetic, I have followed my inclination entirely, without duly considering how little, with my limited powers, I may hope to accomplish anything substantial in this difficult part of mathematics. Although I become more convinced every day of the difficulties connected with investigations of this kind, this occupation has nevertheless become too much a habit for me, and I would almost say too much a passion, for me to decide so easily to abandon the investigations once begun.

Toward the end of this year I shall return to my native country, Prussia, where I hope to be appointed through the kind offices of Baron von \scn{Humboldt}. Although I know how properly to value the advocacy of this great scholar, I cannot help admitting that it still leaves something to be desired. For it depends largely on the judgment that will be passed on my work in Berlin what sphere of activity will be assigned to me in my fatherland, and therefore, in view of the slight value of my attempt, it is of the highest importance to me that it be judged with indulgence, or at least that justice be done to it. But since even among the most distinguished mathematicians -- as I have had occasion here to convince myself -- only very few are familiar with indeterminate analysis, there is reason to fear that my treatise will receive a less favourable reception than might perhaps fall to it if, with otherwise equal intrinsic value, it had as its subject a problem from a better-known part of science, for example astronomy or the integral calculus. This circumstance will excuse me in some measure for daring to trouble Your Excellency with the most submissive request that you might, on some occasion, communicate your judgment of my work to one of the Berlin scholars with whom you are in correspondence. I believe I may entertain the hope that Your Excellency will not refuse this favour to a young German whose fortune would be so greatly advanced by your kind advocacy.

Since I intend to remain in Paris for some time, I take the liberty of offering you my services during my stay in this city, while at the same time begging Your Excellency, should you wish to honour me with a commission or a few lines, kindly to address your letter to my father, Post Commissioner in Düren near Aachen.

\begin{flushright}
I remain in deepest veneration,\\[0.4em]
Your Excellency's\\
most obedient\\
\scn{G. Lejeune Dirichlet}
\end{flushright}
Paris, 28 May 1826.

\section*{II. Gauss to Lejeune Dirichlet}
\begin{flushleft}
To Monsieur\\
\hspace*{2em}Monsieur \scn{Lejeune Dirichlet} in Paris.
\end{flushleft}

I would already have expressed to you earlier my thanks for the treatise you so kindly sent me, and for the great pleasure you thereby gave me, had I not wished first to learn something of the success of what I have tried to do in Berlin with regard to your wishes, and I may add my own. I am extremely pleased now to see from a letter received from the Secretary of the Academy in Berlin that we may hope people will soon be inclined to provide you with an appropriate position in your fatherland.

It is for me all the more pleasing a phenomenon that you devote yourself with great inclination to that part of mathematics which has always been my favourite study, the rarer such an inclination is. I heartily wish you an external position in which, as much as possible, you remain master of your time and of the choice of your work. I myself, immediately after the appearance of my \emph{Disquisitiones}, was very much hindered by employments of another kind, and later by my external circumstances, from following my inclination to the extent I would have wished. Instead of a second part of that work, which I had earlier intended, I shall in all probability have to restrict myself to providing, from time to time, a memoir on a single subject. The three papers of this kind which have so far appeared in the sixteenth volume of the local Commentationes and in the first and fourth volumes of the \emph{Commentationes recentiores} contain, however (apart from part of the second), none of the subjects which I already had in view in 1801 for a continuation, but new ones; and likewise my later works of this kind also concern a new subject, namely the theory of biquadratic residues, which I intend to give in about three papers. The first of these will shortly be printed for the sixth volume of the \emph{Comment. rec.}, and the main materials for the remainder, as well as for the analogous theory of cubic residues, although little of it has yet been put properly on paper, may be regarded as essentially settled.

Kindly commend me to Herr von \scn{Humboldt}, if he is still in Paris, and excuse me for not writing to him himself now, since I am afraid that my letter might not reach him, as I hear he intended to leave Paris.

\begin{flushright}
With sincere esteem,\\
your devoted\\
\scn{C. F. Gauss}
\end{flushright}
Göttingen, 13 September 1826.

\section*{III. Lejeune Dirichlet to Gauss}
\begin{center}\emph{Most highly revered Court Councillor!}\end{center}

I would already have made use of the permission so kindly granted to me by Your Excellency to write to you occasionally, in order to thank you for the goodwill with which you received me, had I not wished at the same time to send you a work with which I began to occupy myself soon after my arrival in Breslau. This work was prompted by the announcement in your notices of your paper on biquadratic residues. The elegance of the criteria communicated in that announcement for deciding whether $\pm2$ is a biquadratic residue or non-residue of a prime number $=8n+1$, the second of which seemed to me especially remarkable, aroused in me the wish to find proofs of them myself. As soon as I had become somewhat familiar here with my professional duties, I began to occupy myself seriously with this subject. My efforts remained fruitless for some time, until I succeeded in deriving your second criterion, whose proof, according to the course chosen by you, seems to require many auxiliary investigations, from the first. But after this first happy success my work again came to a serious halt, and for a long time I could find no means suitable for establishing the first criterion. Finally, toward the beginning of winter, I came upon the proof contained in the accompanying paper, which is so simple that one can hardly understand how it does not at once present itself, once one is familiar with the theorem to be proved. The continuation of my investigations, of which my paper contains only one part, has led me to a great number of results which one certainly would not have suspected in advance to stand in any relation to this subject, and in this work several remarkable examples of the often wonderful connection among arithmetical truths have presented themselves to me, in which you find the principal ground of the attraction which investigations of indeterminate analysis hold for us.

Since I could assume that the readers of my paper would already be somewhat familiar with indeterminate analysis, I expressed myself very briefly in the preliminaries. This circumstance has occasioned a restriction harmful to the generality of the investigation. By analogy I allowed myself to be led into the error that in the theory of biquadratic residues it is sufficient to determine the behaviour of any one prime number with respect to another prime number of the form $4n+1$, whereas the complete answer to all the questions belonging here requires that one also determine the behaviour of composite numbers, at least of those that are products of two primes, with respect to primes of the form $4n+1$. Thus, for example, the question whether $21$ is a biquadratic residue or non-residue of a prime number $84n+1$, $25$, $37$ may very easily be answered by means of the theorems set up in the paper; but the case is different for primes $p$ of the form $84n+5$, $17$, $41$, which require new criteria. For the special case just mentioned I find the following criterion.

\begin{quote}
Put $5p=\varphi^2+\psi^2$ (where I assume that $\psi$ is divisible by $6$). Then $21$ will be a biquadratic residue of $p=84n+5$, $17$, $41$ if one of the numbers $\varphi$, $\psi$ is divisible by $7$; in the contrary case $21$ is a biquadratic non-residue of $p$.
\end{quote}

I am now busily occupied in filling the gap in my work produced by the error indicated above.

During my stay in Göttingen you had the kindness to communicate to me some of the remarkable results that are the fruit of your recent geometrical researches. The wish, already arising then, to become acquainted with your investigations has become very lively in me since, through the abstract of your paper contained in the notices, I have obtained a clearer and more complete idea of the results only lightly indicated in conversation. On this occasion I have again felt very strongly how desirable it would be for learned societies, when publishing their Commentaries, to follow the example given by the London Society, which has a volume appear every half-year. Even your paper on biquadratic residues has still not reached me, although I have repeatedly urged the bookseller to procure it quickly, and the volume of your Commentaries which is to contain it was already announced in the previous catalogue.

May I ask Your Excellency to accept the assurance of the heartfelt veneration and gratitude with which I remain

\begin{flushright}
Your Excellency's\\
obedient\\
\scn{Lejeune Dirichlet}
\end{flushright}
Breslau, 8 April 1828.

\section*{IV. Gauss to Lejeune Dirichlet}

For your kind letter and the agreeable sending of your two papers, I offer you, my highly esteemed friend, my most binding thanks. I see with pleasure the increasing interest that is now beginning to be taken in investigations of higher arithmetic. The happy way in which you derive the second criterion relating to the biquadratic residuacity of the number $2$ from the first pleased me very much.

Presumably the sixth volume of our Commentationes has now found its way to Breslau, and my \emph{Commentatio prima} on biquadratic residues is therefore probably known to you. Should an opportunity present itself, I would also gladly send you a separate offprint of it. I could have chosen among several kinds of proof for the theorem occurring there; but it will not escape you why I preferred the one carried out there, chiefly because the classification of $2$ for those moduli for which it is a quadratic non-residue (under $B$ or $D$) must be regarded as an essential integral part of the theorem, to which most of the other methods of proof do not seem applicable.

The whole investigation, the material of which I have possessed completely for 23 years, but the proofs of the main theorems (to which the theorem in the first Commentation is not yet to be counted) for about 14 years -- although I wish and hope still to simplify something in the latter, the proofs -- I have calculated at about three papers. I have already now begun drafting the second, and hope to complete it in not too long a time, unless the surveying business recently imposed on me again causes some delay.

The final theorem $b\equiv\frac12 rr'\pmod p$ I had already made known three years ago in the local learned notices, and I had called attention to the remarkable knot that remained to be solved there. But so far I have not heard that anyone attempted it. A few days ago I really succeeded with one half, and this discovery gave me all the more pleasure because it is not at all founded on induction -- for I admit that I would not have expected precisely this connection, but rather a priori from the combination of other, very entangled and interesting investigations, already 28 years old but not yet made known at all, of which a slight indication is given in the concluding note to the \emph{Disquis. Arithm.} p. 668.

This is namely a sufficient criterion for the case in which $p$ is of the form $8n+5$.

Let the number of classes which the binary forms in each of the two kinds form for determinant $-p$ be equal to $k$. The beginning of a table I constructed up to determinant $-3000$ stands in the \emph{Disquis. Arithm.} p. 520. It is also still to be noted that, for a $p$ of the assumed form, one always has $k=2m+1$ when $m$ denotes the number of decompositions of $p$ into three positive squares (I say positive in order to exclude $0$), as \scn{Legendre} found by induction and as was first proved in the \emph{Disquisitiones Arithm.} from the theory of ternary forms. For example:

\begin{center}\scriptsize
\begin{tabular}{c|cccccccccc}
$p$&5&13&29&37&53&61&101&109&149&157\\
$k$&1&1&3&1&3&3&7&3&7&3\\
$m$&0&0&1&0&1&1&3&1&3&1\\
decompositions&&&$4+9+16$&&$1+16+36$&$9+16+36$&\begin{tabular}{c}$1+36+64$\\$4+16+81$\\$16+36+49$\end{tabular}&$9+36+64$&\begin{tabular}{c}$1+4+144$\\$4+64+81$\\$36+64+49$\end{tabular}&$4+9+144$
\end{tabular}\end{center}

Assuming this, that value of $b$ which is congruent to $\frac12 rr'\pmod p$ is always
\[
 b\equiv2k+a-1\equiv4m+a+1\pmod 8,
\]
by which the sign of $b$ is completely determined. Here are 22 examples, by which I double the extent of the table given at the end of the paper:

\begin{center}\small
\begin{tabular}{rrrr|rrrr}
\toprule
$p$&$k$&$a$&$b$&$p$&$k$&$a$&$b$\\
\midrule
5&1&$+1$&$+2$&181&5&$+9$&$+10$\\
13&1&$-3$&$-2$&197&5&$+1$&$-14$\\
29&3&$+5$&$+2$&229&5&$-15$&$+2$\\
37&1&$+1$&$-6$&269&11&$+13$&$+10$\\
53&3&$-7$&$-2$&277&3&$+9$&$+14$\\
61&3&$+5$&$-6$&293&9&$+17$&$+2$\\
101&7&$+1$&$-10$&317&5&$-11$&$+14$\\
109&3&$-3$&$+10$&349&7&$+5$&$+18$\\
149&7&$-7$&$-10$&373&5&$-7$&$+18$\\
157&3&$-11$&$-6$&389&11&$+17$&$-10$\\
173&7&$+13$&$+2$&397&3&$-19$&$-6$\\
\bottomrule
\end{tabular}\end{center}

One can therefore also express the rule as follows, always assuming $p\equiv5\pmod8$. One has
\[
 b\equiv a+1\pmod8\quad\text{if }m\text{ is even},
\]
and
\[
 b\equiv a+5\pmod8\quad\text{if }m\text{ is odd}.
\]
I do not yet venture a conjecture whether an even simpler criterion is possible, by which one could distinguish in advance the case of even $m$ from that of odd $m$, that is, without knowing the value of $m$ itself, since, as I have noted above, this rapprochement is still quite new. In the case $p\equiv1\pmod8$, the above congruence $b\equiv2k+a-1\pmod8$ does remain correct, but it no longer decides the sign of $b$, since it is simultaneously satisfied by the positive and the negative value of $b$. Here $k$ is namely always even, $=2m$ (if the meaning of $m$ is stated as above) or $=2m+2$ if by $m$ one understands the number of decompositions of $p$ into three positive unequal squares, and $b\equiv0\pmod4$, or $b\equiv -b\pmod8$. I suppose that the case $p\equiv1\pmod8$, or $b\equiv0\pmod4$, is \emph{altioris indaginis} and perhaps again
\[
 b\equiv4\pmod8 \quad \text{is easier than}\quad b\equiv0\pmod8,
\]
\[
 b\equiv8\pmod{16}\quad \text{is easier than}\quad b\equiv0\pmod{16},\quad\text{etc.}
\]

\begin{flushright}
With distinguished esteem I remain\\
your friendly devoted\\
\scn{C. F. Gauss}
\end{flushright}
Göttingen, 30 May 1828.

\section*{V. Lejeune Dirichlet to Gauss}

Receive, highly honoured patron, my most heartfelt thanks for the permission so kindly granted to me to spend some time near you. Unfortunately I am not able to make use of it during these holidays and to carry out a long-cherished wish. A few days ago I had a rather violent attack of dysentery, and now, when I am beginning to recover, my family are laid low with measles. I need hardly tell you how much I regret seeing my visit to Göttingen postponed once again. Occupied for years with the study of your works, nothing could be more enjoyable and more advancing for me than to owe to your kindness explanations of the origin and development of the splendid methods which appear in your writings in admirable perfection. Indeed, I will confess to you that in my hopes I go so far as to expect, from your goodwill, oral communication of part of your theoretical investigations on magnetism and electromagnetism, and I solemnly vow in advance, should you deem me worthy of such a favour, to take Herr \scn{Krone}'s discretion as my model; he, in the enviable possession of this treasure, would reveal nothing except that these investigations are intimately connected with the beautiful method developed in your paper on the attraction of ellipsoids. Have you learned from the \emph{Comptes rendus} of the Paris Academy of the discussion which has recently arisen over the problem just mentioned? \scn{Poisson} seems to regard his solution as the only direct one, while in all others, in the case of an exterior point, the difficulty is merely eluded. I must frankly confess that statements of this kind have no clear sense for me. To reduce a double integral to a simple one by introducing new variables for which one integration can be performed does not seem to me more direct than any other means, for example differentiation under the integral sign. In my opinion everything depends on the route by which the reduction can be carried out most quickly and clearly; and in this respect your treatment, which gives the attraction for interior and exterior points at one stroke, has decisively the advantage over all known ones. If such distinctions were admitted, then, to be consistent, one would have to admit for almost all definite integrals that in their discovery the difficulty was merely eluded; yet by this, it seems to me, actually nothing at all would be said.

Permit me, highly honoured patron, to send you, enclosed, a short article which appeared in the last issue of Crelle's Journal. I was led to the method indicated there, which appears to me very fruitful, through the effort to prove the theorem that every arithmetic progression whose first term and difference have no common factor contains infinitely many primes. I would almost conjecture that my method has some analogy with the investigations indicated in the final note of the \emph{Disq. arith.}, especially because you say of your investigations that they shed light on several parts of analysis, while my method stands in such intimate connection with the remarkable trigonometric series which represent discontinuous functions and whose nature, at the time of the appearance of the \emph{Disq. arith.}, was still wholly unexplained. You may perhaps still remember that more than ten years ago, when I was still in Breslau, you communicated to me a criterion for deciding the question raised at the end of your first paper on biquadratic residues (for the case $p=8n+5$), with the remark that you had derived it from the investigations announced at the end of the \emph{Disq. arith.} This criterion now follows very easily from the combination of the theorem that for determinant $-p$ the number of forms in each genus is equal to the excess of the number of quadratic residues $<\frac p4$ over the number of quadratic non-residues $<\frac p4$, with the theory of biquadratic residues, this theory itself taken only to the extent in which I have given it in Crelle's Journal. In a detailed working out of my investigations on the use of series in higher arithmetic, with which I intend shortly to occupy myself, may I mention this criterion as having been communicated to me by you in 1828?

Allow me to ask you to hand the enclosed letter to \scn{Weber}, and kindly accept the assurance of my deepest admiration and heartfelt veneration.

\begin{flushright}
Yours gratefully devoted,\\
\scn{Dirichlet}
\end{flushright}
Berlin, 9 September 1838.

\section*{VI. Lejeune Dirichlet to Gauss}
\begin{center}\emph{Most revered patron!}\end{center}

The low spirits caused in me by the illness from which my wife has recently suffered will, I hope, excuse me for having delayed so long in writing to you, however much I wished to offer you my heartfelt thanks for the exceedingly kind and benevolent reception you granted me. Another reason for postponing my intention lay in the fact that several of the books I wanted to consult, in order to give you the desired information about what has been written so far on the shortest surface, were just not in the library, and I therefore had to await their return.

I have found nothing at all on the subject named, but I have found investigations on other problems which, although different in object, coincide in essence with the determination of the shortest surface. These problems concern the theory of heat in the case of permanent temperatures, to which, according as one takes two or three dimensions into account, the equation
\[
(1)\quad \frac{\partial^2u}{\partial x^2}+\frac{\partial^2u}{\partial y^2}=0,
\qquad\text{or}\qquad
(2)\quad \frac{\partial^2u}{\partial x^2}+\frac{\partial^2u}{\partial y^2}+\frac{\partial^2u}{\partial z^2}=0
\]
relates. If the given curve, only slightly removed from a plane, through which the shortest surface is to be laid, projects as a rectangle or a circle, then the solution of the problem is easily derived from the methods given by \scn{Fourier} in his \emph{Théorie de la Chaleur}. The much more difficult case in which the curve projects as an ellipse also seems to be contained implicitly in a beautiful paper by \scn{Lamé} (Journal de Liouville, May 1839), since there a function $u$ is determined which, inside an ellipsoid given by
\[
 \left(\frac{x}{\alpha}\right)^2+\left(\frac{y}{\beta}\right)^2+\left(\frac{z}{\gamma}\right)^2=1,
\]
satisfies differential equation (2) and on the surface reduces to an arbitrary function of the two coordinates determining each point of it. If one sets $\gamma=\infty$ and lets the given arbitrary function become independent of $z$, one obtains the equation of the shortest surface whose boundary curve projects as an ellipse.

Since my return I have occupied myself somewhat with quadratic forms for complex numbers and have convinced myself that almost all the results and methods of section 5 of the \emph{Disq. arith.} remain applicable to this case \emph{mutatis mutandis}. In a similar way to the case of real numbers, the number of forms can also be determined; the remarkable result arises that this number, as it is connected for real positive determinants with the division of the circle, stands here in an analogous relation to the division of the lemniscate. These investigations have led me to a theorem which seems not unimportant for the theory of indeterminate equations of higher degree and which offers a remarkable generalization of the theorem according to which the equation $t^2-Du^2=1$ is always solvable.

Let $a,b,\ldots,g,h$ be integers, and suppose the equation
\[
 s^n+as^{n-1}+bs^{n-2}+\cdots+gs+h=0
\]
has no rational factor and has among its roots $\alpha,\beta,\gamma,
\ldots$ at least one real root. Then integers
\[
 \lambda,\mu,\nu,\ldots,\varrho\qquad(\text{not with }\lambda=1,\ \mu=0,\ \nu=0,\ldots)
\]
can always be found, and in infinitely many different ways, of such a kind that, if for abbreviation one puts
\[
 \lambda+\mu\alpha+\nu\alpha^2+\cdots+\varrho\alpha^{n-1}=F(\alpha),
\]
the product
\[
 F(\alpha)F(\beta)F(\gamma)\cdots=1
\]
holds. The proof of this theorem is in the highest degree simple and can be dispatched with the aid of geometrical considerations in two or three pages.

I cannot close this letter without repeating the wish already expressed to you in Göttingen, that you might be pleased to favour me with a separate offprint of your beautiful investigations on the general theory of attractions. I would most reluctantly have to wait until the completion of the whole next issue of the \emph{Resultate} etc. to become acquainted with them. In expectation of your kind granting of my request, I remain

\begin{flushright}
your gratefully devoted\\
\scn{Lejeune Dirichlet}
\end{flushright}
Berlin, 3 January 1840.

May I ask you to hand the enclosed letter to \scn{Weber}, or, if he should be in Leipzig, to send it there by post.

\section*{VII. Lejeune Dirichlet to Gauss}
\begin{center}\emph{Highly honoured patron!}\end{center}

If I communicate to you so late the result of the inquiries undertaken at your commission concerning the Russian books, this delay, I hope, will find some excuse in the good success of my efforts. Like you, Privy Councillor von \scn{Varnhagen} found great difficulties in obtaining Russian books. The connections formerly made by him with great effort for this purpose now largely no longer exist, since for years he has greatly neglected his Russian studies and consequently lacked opportunity to cultivate the earlier connections. In that earlier period, however, when he pursued Russian with great zeal, he acquired a not inconsiderable stock of books, which he places at your disposal for whatever use you wish, with great readiness. At my request he drew up a list of the most important of these books, which I enclose. Please kindly tell me which of these works you wish to read. I shall then immediately ask H. von V. for what you desire and send it to you without loss of time.

I returned from my holiday journey just in time to see poor \scn{Eisenstein} at least once more. When I visited him the day after my arrival, I found him changed to a frightening degree and so wholly occupied with his suffering that even the attempt to tell him of you and so to distract him for a moment remained without success. With such boundless suffering, the death that occurred that same night had to appear to me as a benefit for our poor friend.

As I hear from \scn{Weber}, you desire information on the origin of the mathematical library which is to be sold here in March, unless it is earlier acquired as a whole -- for which there seems to be some likelihood. This library was formerly in the possession of a Privy Councillor \scn{Heiligenstein}, who stayed in Paris at the beginning of the century, and then later passed into the possession of his son, who is said to have been an astronomer. With this library, for simultaneous auctioning, some other things have been joined, among which are, for example, the two manuscripts by you -- of which one is that of the \emph{Determinatio attract.}; I do not immediately recall the content of the other -- which come from \scn{Dirksen}'s estate.

You recently seemed to wish that my proof for the convergence of trigonometric series be extended to the case in which the maxima and minima of the function to be expanded are no longer countable, and at the same time you expressed the supposition that the proof of the theorem under this extended hypothesis would probably offer no essential difficulties. On closer consideration I find your supposition completely confirmed, provided one disregards certain quite singular cases. The whole foundation of the proof is formed by the theorems according to which
\[
 \lim \int_0^c f(\beta)\frac{\sin k\beta}{\sin\beta}\,d\beta=\frac{\pi}{2}f(0),
 \qquad\text{and}\qquad
 \lim \int_b^c f(\beta)\frac{\sin k\beta}{\sin\beta}\,d\beta=0,
\]
where $0<b<c\leq\frac\pi2$ is assumed and the symbol $\lim$ refers to the positive quantity $k$ becoming infinite.

The extension of these theorems to the case just mentioned, in which $f(\beta)$ presents an infinite number of maxima and minima, is quite easy, at least for the second theorem and also for the first, provided only that infinitely many maxima and minima do not lie in the immediate neighbourhood of $\beta=0$. Namely, one need only separate from the integrals those parts within which such an accumulation of maxima and minima occurs, and one can arrange these, by bringing the bounds of these partial integrals sufficiently close to each other, so that they remain arbitrarily small, however large $k$ becomes. From this simple remark, and by the way in which convergence is concluded from the theorems above, it already follows that the trigonometric series for representing an arbitrary function $\varphi(x)$ always converges even when the maxima and minima of $\varphi(x)$ are no longer countable, provided only that one assigns to $x$ no value in the immediate neighbourhood of which an infinite accumulation of maxima and minima of $\varphi(x)$ occurs. To prove convergence for the exceptional values just mentioned, it must be shown that the first of the above theorems remains valid when $f(\beta)$, from $\beta=0$ to $\beta=\delta$, where $\delta$ is arbitrarily small, passes infinitely often from increasing to decreasing and conversely. If then $f(\beta)$ is so constituted that the difference $f(\beta)-f(0)$ can be represented as a product whose first factor $\varphi(\beta)$ remains finite for infinitely small $\beta$, while the second, which I suppose positive and shall call $\psi(\beta)$, has the property that the integral
\[
 \int_0^\delta \psi(\beta)\frac{d\beta}{\beta},
\]
as is the case for example for every positive power $\beta^p$, itself becomes infinitely small when $\delta$ is infinitely small, then one may decompose
\[
 \int_0^\delta f(\beta)\frac{\sin k\beta}{\sin\beta}\,d\beta
\]
only into the two components
\[
 f(0)\int_0^\delta \frac{\sin k\beta}{\sin\beta}\,d\beta
 +\int_0^\delta \varphi(\beta)\sin k\beta\frac{\psi(\beta)}{\sin\beta}\,d\beta.
\]
Of these, for fixed $\delta$, the first passes as $k$ increases into $\frac\pi2 f(0)$, while the second, for a suitably small chosen $\delta$, always remains smaller than an arbitrarily small quantity. From this, and since
\[
 \int_\delta^c f(\beta)\frac{\sin k\beta}{\sin\beta}\,d\beta
\]
has zero as its limit, the validity of the first theorem follows, at least under the earlier hypothesis on $f(\beta)-f(0)$.

Forgive, revered patron, this long, almost unletterlike mathematical exposition, which I would not have dared to write had I foreseen, just as I began it, its great extent.

Please accept, highly honoured patron, the expression of my deep admiration and heartfelt veneration.

\begin{flushright}
yours gratefully\\
devoted\\
\scn{Dirichlet}
\end{flushright}
Berlin, 20 February 1853.

\end{document}
